\section{Local Connectivity Preserving Complete h-Splitting-Off}
\label{sec:hypergraph-splitting-off-proof}
We prove Theorem \ref{thm:complete-splitting-off} using Theorem \ref{thm:WeakToStrongCover:main} in this section. 
We need the following lemma showing that a certain set function $p: 2^V\rightarrow \Z$ is symmetric skew-supermodular and admits an efficient algorithm to implement \functionMaximizationOracleStrongCover{p}. This lemma has been observed in the literature before. We refer the reader to \Cref{sec:helper-for-p-max-oracle} for its proof. 

\begin{restatable}{lemma}{lempmaxoracle}\label{lemma:helper-for-p-max-oracle}
    Let $(G=(V, E), c: E\rightarrow \Z_+)$ be a hypergraph. Let $r:\binom{V}{2}\rightarrow \Z_{\ge 0}$ be a function on pairs of elements of $V$ and $R, p_{(G, c, r)}:2^V\rightarrow \Z$ be functions defined by $R(X)\coloneqq \max\{r(\{u, v\}): u\in X, v\in V\setminus X\}$ for every $X\subseteq V$, $R(\emptyset)\coloneqq 0$, $R(V)\coloneqq 0$, and $p_{(G, c, r)}(X)\coloneqq R(X)-d_{(G, c)}(X)$ for every $X\subseteq V$. 
    \begin{enumerate}
        \item The function $p_{(G, c, r)}$ is symmetric skew-supermodular. 
        \item 
        Given hypergraph $(G, c)$, function $r$, hypergraph $(G_0=(V, E_0), c_0: E_0\rightarrow \Z_{+})$, and disjoint sets $S_0, T_0\subseteq V$, the oracle   \functionMaximizationOracleStrongCover{p_{(G, c, r)}}$((G_0, c_0), S_0, T_0)$ can be computed in $O(|V|^3(|V|+|E_0|+|E|)(|E_0|+|E|))$ time.
        \iffalse
        Given $(G, c)$ and the function $r$, there exists a strongly polynomial-time algorithm to implement \linebreak \functionMaximizationOracleStrongCover{p}. 
        In particular, for a given hypergraph $(G_0=(V, E_0), c_0: E_0\rightarrow \Z_{+})$, and disjoint sets $S_0, T_0\subseteq V$, the oracle   \functionMaximizationOracleStrongCover{p}$((G_0, c_0), S_0, T_0)$ can be implemented to run in $O(|V|^3(|V|+|E_0|+|E|)(|E_0|+|E|))$ time.
        \fi
    \end{enumerate}
\end{restatable}


We now prove the following stronger form of Theorem \ref{thm:complete-splitting-off}.
\begin{theorem}\label{thm:hypergraphs-splitting-off-stronger-form}
    Let $(G=(V, E), w_G: E\rightarrow \Z_+)$ be a hypergraph and $s\in V$. Then, there exists a hypergraph $(G^*=(V, E^*), c^*)$ such that 
    \begin{enumerate}
        \item $(G^*, c^*)$ is a local connectivity preserving complete h-splitting-off at $s$ from $(G, w_G)$ and 
        \item $|E^*|-|E|=O(|V|)$. 
    \end{enumerate}
    Moreover, given the hypergraph $(G, w_G)$, a hypergraph $(G^*, c^*)$ satisfying the above two properties can be obtained in $O(|V|^6(|V|+|E|)^3)$ time. 
\end{theorem}
\begin{proof}
    We begin by proving the existence of a local connectivity preserving complete h-splitting-off at $s$ from $(G=(V, E), w_G: E\rightarrow \Z_+)$ such that the number of additional hyperedges is $O(|V|)$. 
    We will use \Cref{thm:WeakToStrongCover:main} with an appropriate setting of the symmetric skew-supermodular function $p$ and the hypergraph $(H, w)$. 
    Let $V'\coloneqq V\setminus \{s\}$. 
    Let the hypergraph $(H:=(V', E_H), w_H: E_H\rightarrow \Z_+)$ be defined by $E_H\coloneqq \{e-\{s\}: e\in \delta_G(s)\}$ with $w_H(e-\{s\})\coloneqq w_G(e)$ for every $e\in \delta_G(s)$. 
    In order to define the function $p: 2^{V'}\rightarrow \Z$, we consider another hypergraph $(J:=({V'}, E_J), w_J: E_J\rightarrow \Z_+)$ defined by $E_J\coloneqq E\setminus B_G(s)$ with $w_J(e)\coloneqq w_G(e)$ for every $e\in E\setminus B_G(s)$. 
    %We recall that the cut function $d_{(J, w_J)}: 2^{V'}\rightarrow \Z_{\ge 0}$ of the hypergraph $(J, w_J)$ is symmetric submodular. 
    We consider the function $R:2^{V'}\rightarrow \Z_{\ge 0}$ defined by $R(X)\coloneqq \max\{\lambda_{(G, w_G)}(u, v): u\in X, v\in V'\setminus X\}$ for every non-empty proper subset $X\subsetneq  V'$, $R(\emptyset)\coloneqq 0$, and $R(V')\coloneqq 0$. 
    %We observe that the function $R$ is symmetric skew-supermodular \knote{need to prove or cite a reference}. 
    Now, consider the function $p: 2^{V'}\rightarrow \Z$ defined by $p(X)\coloneqq R(X)-d_{(J, w_J)}(X)$ for every $X\subseteq {V'}$. 
    %Since $p$ is the difference of a symmetric skew-supermodular function and a symmetric submodular function, it follows that 
    The function $p$ is symmetric skew-supermodular by \Cref{lemma:helper-for-p-max-oracle}. 
    We also note that $b_{(H, w_H)}(X)\ge p(X)$ for every $X\subseteq {V'}$: for an arbitrary $X\subseteq {V'}$, we have that $b_{(H, w_H)}(X)+d_{(J, w_J)}(X)=d_{(G, w_G)}(X)\ge R(X)$ where the last inequality is by definition of the function $R$ and hence, $b_{(H, w_H)}(X)\ge R(X)-d_{(J, w_J)}(X)=p(X)$. 
    Hence, $(H,w_H)$ weakly covers the symmetric skew-supermodular function $p$. 
    Applying \Cref{thm:WeakToStrongCover:main} to the hypergraph $(H, w_H)$ and the function $p$, we obtain that there exists a hypergraph $(H^* = (V', E_{H^*}), w_{H^*}: E_{H^*}\rightarrow \Z_+)$ satisfying the following three properties: (i) $d_{(H^*, w_{H^*})}\ge p(X)$ for every $X\subseteq V'$,  (ii) the hypergraph $(H^*, w_{H^*})$ is obtained by merging hyperedges of the hypergraph $(H, w_H)$, and (iii) $|E_{H^*}|-|E_H|=O(|V|)$. 
    Consider the hypergraph $(G^*=(V, E^*), c^*)$ obtained from the hypergraph $(J+H^*, w_J+w_{H^*})$ by adding the vertex $s$. 
    We have that 
    \begin{align*}
    |E^*|
    &=|E_J\cup E_{H^*}| &\text{(by definition of $E^*$)}\\
    &\le |E_J|+|E_{H^*}| &\\
    &=|E_J|+|E_H|+O(|V|) & \text{(by property (iii))}\\
    &=|E_J\cup E_H|+O(|V|) & \text{(by definitions of $E_J$ and $E_H$)}\\
    &=|E|+O(|V|), & \text{(by definitions of $E_J$ and $E_H$)}
    \end{align*}
    and hence, $|E^*|-|E|=O(|V|)$. 
    
    We now show that $(G^*, c^*)$ is a complete h-splitting-off at $s$ from $(G, w_G)$ and $(G^*, c^*)$ preserves local connectivity. 
    We note that property (ii) implies that $(G^*, c^*)$ is a complete h-splitting-off at $s$ from $(G, c)$: this is because, 
    hyperedges of $E_{H^*}\setminus E_H$ that are obtained by merging hyperedges of $E_H$ are equivalent to being obtained by a sequence of h-merge almost-disjoint hyperedges operation starting from the hypergraph $(G, w_G)$ and a final h-trim operation; 
    moreover, hyperedges of $E_{H^*}\cap E_H$ are equivalent to being obtained by the h-trim operation from the hypergraph $(G, w_G)$;  
    finally, we have that $d_{(G^*, c^*)}(s)=0$ by definition of $(G^*, c^*)$. 
    Next, we note that property (i) implies that $(G^*, c^*)$ preserves local connectivity: consider an arbitrary pair of distinct vertices $u, v\in V'$. We recall that h-merge almost-disjoint hyperedges as well as h-trim operations do not increase the cut value of any subset $X\subseteq V'$ and hence, $\lambda_{(G^*, c^*)}(u, v) \le d_{(G^*, c^*)}(X)\le d_{(G, w_G)}(X)$ for every $X\subseteq V'$ such that $u\in X, v\in V'\setminus X$ and in particular, $\lambda_{(G^*, c^*)}(u, v) \le \lambda_{(G, w_G)}(u, v)$. Hence, it suffices to show that $d_{(G^*, c^*)}(X)\ge \lambda_{(G, w_G)}(u, v)$ for every $X\subseteq V'$ such that $u\in X, v\in V'\setminus X$. Let $X\subseteq V'$ such that $u\in X, v\in V'\setminus X$. We have that 
    \begin{align*}
    d_{(G^*, c^*)}(X)
    &=d_{(J, w_J)}(X) + d_{(H^*, w_{H^*})}(X) & \text{(by definition of $(G^*, c^*)$)}\\
    &\ge d_{(J, w_J)}(X) +p(X) & \text{(by property (i) of $(H^*, w_{H^*})$)}\\
    &=d_{(J, w_J)}(X) +R(X)-d_{(J, w_J)}(X) & \text{(by definition of the function $p$)}\\
    &=R(X) \\
    &\ge \lambda_{(G, w_G)}(u,v). & \text{(by definition of the function $R$)}
    \end{align*}
    %Finally, we note that $b_{(G^*, c^*)}(s)=0$ by definition. 
    This completes the proof that $(G^*, c^*)$ is a complete splitting off at $s$ from $(G, w_G)$ and preserves local connectivity. 
    
    %We now show a strongly polynomial-time algorithm to obtain such a hypergraph $(G^*, c^*)$. Given $(G, w_G)$ as input, we can construct $(H, w_H)$ and $(J, w_J)$ as above in strongly polynomial time. There exists a strongly polynomial-time algorithm to implement \functionMaximizationOracleStrongCover{p}, where $p$ is the function defined above (see \Cref{lemma:helper-for-p-max-oracle}). By the algorithmic conclusion of \Cref{thm:WeakToStrongCover:main}, it follows that the hypergraph $(H^*, w_{H^*})$ mentioned in the previous paragraph can be obtained in strongly polynomial time. Finally, we can construct the hypergraph $(G^*, c^*)$ using the hypergraph $(H^*, w_{H^*})$ and the hypergraph $(J, w_J)$ in strongly polynomial time. 
   We now analyze the run-time of the algorithm to obtain such a hypergraph $(G^*, c^*)$. Given $(G=(V, E), w_G: E\rightarrow \Z_+)$ as input, we can construct $(H, w_H)$ and $(J, w_G)$ as above in time $O(|V| + |E|)$. We note that both $(H, w_H)$ and $(J, w_G)$ have vertex set $V-\{s\}$ and at most $|E|$ hyperedges. For the function $p$ defined as above, we consider the time to implement \functionMaximizationOracleStrongCover{p}: by Lemma \ref{lemma:helper-for-p-max-oracle}, for a given hypergraph $(G_0=(V-\{s\}, E_0), c_0: E_0\rightarrow \Z_+)$, sets $S_0, T_0\subseteq V$, and a vector $y_0\in \R^V$, the oracle \functionMaximizationOracleStrongCover{p}$((G_0, c_0), S_0, T_0, y_0)$ can be implemented to run in $O(|V|^3(|V|+|E_0|+|E|)(|E_0|+|E|))$ time. By the algorithmic conclusion of \Cref{thm:WeakToStrongCover:main}, it follows that the hypergraph $(H^*, w_{H^*})$ mentioned in the previous paragraph can be obtained in $O(|V|^6(|V|+|E|)^3)$ time. Finally, we can construct the hypergraph $(G^*, c^*)$ using the hypergraph $(H^*, w_{H^*})$ and the hypergraph $(J, w_J)$ in $O(|V|+|E_{H^*}|+|E|)=O(|V|+|E|)$ time. Thus, the overall run-time is $O(|V|^6(|V|+|E|)^3)$.
\end{proof}
